% --- IMPORTS ---
\documentclass[leqno]{article}
\usepackage{graphicx}
\usepackage{dsfont}
\usepackage{mathtools}
\usepackage{amssymb}
\usepackage{amsmath}
\usepackage{amsthm}
\usepackage{float}
\usepackage{ragged2e}
\usepackage{tensor}
\usepackage{fancyhdr}
\usepackage{empheq}
\usepackage[most]{tcolorbox}
\usepackage{enumitem}
\usepackage{dirtytalk}
\usepackage{marginnote}
\usepackage{microtype}
\usepackage[
  a4paper,
  left=0.8in,
  right=1in,
  top=1in,
  bottom=0.5in,
  headheight=14pt,
  headsep=0.4in
]{geometry}
\setlength{\parindent}{0cm}
% --- TikZ Packages ---
\usepackage{pgfplots}
\pgfplotsset{compat=1.18}
\usepgfplotslibrary{fillbetween, polar}
\usetikzlibrary{calc, positioning}
\usetikzlibrary{angles, quotes}
\usetikzlibrary{arrows.meta}
\usetikzlibrary{decorations.pathreplacing, decorations.markings}
\usetikzlibrary{patterns}
\usetikzlibrary{intersections}
% --- URL Packages ---
\usepackage[colorlinks=true,linkcolor=red,citecolor=magenta,urlcolor=black]{hyperref}%
\urlstyle{same}
% --- END OF IMPORTS ---

% --- CUSTOM COMMANDS ---
\RenewDocumentCommand{\marks}{o m}{%
  \IfValueTF{#1}%
    {\marginnote{\raggedright\textbf{[#2]}}[#1]}%
    {\marginnote{\raggedright\textbf{[#2]}}}%
}
\newcommand{\vspacerule}[2]{%
  \par\vspace*{#1}%
  \noindent\hrulefill\par
  \vspace*{#2}%
}
\definecolor{scarlet}{RGB}{205, 40, 75}
\newcommand{\questionitem}{%
  \item\phantomsection
  \addcontentsline{toc}{section}{Question \number\value{enumi}}%
}
\renewcommand{\contentsname}{Questions}
% --- END OF CUSTOM COMMANDS ---

% --- HEADER CONFIG ---
\pagestyle{fancy}
\fancyhead{}
\fancyfoot{}
\renewcommand{\headrulewidth}{0.9pt}
\lhead{Geometry of Plane Intersections}
\chead{}
\rhead{\href{https://danieljsmith.org}{danieljsmith.org}}
% --- END OF HEADER CONFIG ---

\begin{document}
\vspace*{20pt}
\tableofcontents
\newpage
\begin{enumerate}[
  label=\textbf{\arabic*.},
  leftmargin=*,
  topsep=0pt,
  partopsep=0pt,
  parsep=0pt
]

% ---- Question 1 ----
\questionitem

Three planes have equations
\[
\begin{aligned}
x-y+2z&=3\\
2x+y+z&=4\\
3x+(m+1)z&=7
\end{aligned}
\] \\
where $m$ is a constant. \\
\begin{enumerate}[label=(\roman*)]
  \item Investigate the arrangement of the planes:

  when $m=2$;

  when $m\ne 2$. \marks{6} \\
  \item A student claims that the position vectors $\mathbf{i}-\mathbf{j}+2\mathbf{k}$, $2\mathbf{i}+\mathbf{j}+\mathbf{k}$ and $3\mathbf{i}+4\mathbf{k}$ all lie in a plane through the origin. Determine whether or not the student is correct. \marks{2} \\
\end{enumerate}

\newpage

% ---- Question 2 ----
\questionitem

\begin{enumerate}[label=\textbf{(\alph*)}]
  \item Consider the linear system $A\mathbf{x}=\mathbf{b}$, where
  \[
  A=\begin{pmatrix}
  3&-2&1\\
  2&1&-4\\
  7&0&-7
  \end{pmatrix},
  \qquad
  \mathbf{x}=\begin{pmatrix}x\\y\\z\end{pmatrix},
  \qquad
  \mathbf{b}=\begin{pmatrix}5\\1\\7\end{pmatrix}
  \]
  Show that $A$ is singular, and deduce that the system does not have a unique solution. \marks{2}\\
  
  \item Show that the system is consistent. Interpret the solution set geometrically. \marks{3} \\
\end{enumerate}

\newpage

% ---- Question 3 ----
\questionitem

\begin{enumerate}[label=\textbf{(\alph*)}]
  \item Show that the system of equations
  \[
  \begin{aligned}
  x+2y-z&=3\\
  2x-y+4z&=1\\
  3x+y+3z&=a
  \end{aligned}
  \]
  where $a$ is a constant, does not have a unique solution. \marks{2}\\
  \item Given that $a=4$, show that the system of equations in part (a) is consistent. Interpret this situation geometrically. \marks{3} \\
  \item Given instead that $a\ne 4$, show that the system of equations in part (a) is inconsistent. Interpret this situation geometrically. \marks{2} \\
\end{enumerate}

\newpage

% ---- Question 4 ----
\questionitem

The equations
\[
\begin{aligned}
x-y+2z&=4\\
(k+1)x+ky+(2-k)z&=4-k\\
(5-2k)x+(2-2k)y+(k+2)z&=8-k
\end{aligned}
\] \\
represent three planes, where $k$ is a real constant. \\
The planes do not intersect in exactly one point. \\
\begin{enumerate}[label=\textbf{(\alph*)}]
  \item Find the possible values of $k$. \marks{3} \\
  \item For each value of $k$ found in part (a), describe the configuration of the planes. \\
  Fully justify your answer, stating in each case whether the corresponding system is consistent or inconsistent. \marks{5} \\
\end{enumerate}

\newpage

% ---- Question 5 ----
\questionitem

Three planes have equations
\[
\begin{aligned}
x+ay&=-1\\
2x+y+z&=3\\
3x+by+2z&=c
\end{aligned}
\] \\
where $a$, $b$ and $c$ are constants. \\
\begin{enumerate}[label=\textbf{(\alph*)}]
  \item In the case where the planes \textbf{do not} intersect at a unique point,
  \begin{enumerate}[label=(\roman*)]
    \item find $b$ in terms of $a$ \marks{4}
    \item find the value of $c$ for which the planes form a sheaf. \marks{3} \\
  \end{enumerate}
  \item In the case where $b=1-a$ and $c=5$, find the coordinates of the point of intersection of the planes in terms of $a$. \marks{6} \\
\end{enumerate}

\newpage

% ---- Question 6 ----
\questionitem

\[
\mathbf{M}=\begin{pmatrix}
1&1&0\\
3&k&1\\
2&-1&1
\end{pmatrix}
\] \\
where $k$ is a constant. \\
\begin{enumerate}[label=\textbf{(\alph*)}]
  \item Find the values of $k$ for which the matrix $\mathbf{M}$ has an inverse. \marks{2} \\
  \item Find, in terms of $p$, the coordinates of the point where the following planes intersect:
  \[
  \begin{aligned}
  x+y&=1\\
  3x+2y+z&=p\\
  2x-y+z&=4
  \end{aligned}
  \]
  \marks[-20pt]{5}
  \item
  \begin{enumerate}[label=(\roman*)]
    \item Find the value of $q$ for which the set of simultaneous equations
    \[
    \begin{aligned}
    x+y&=1\\
    3x+z&=q\\
    2x-y+z&=4
    \end{aligned}
    \]
    has at least one solution. \marks[-20pt]{2}
    \item For this value of $q$, interpret the solution set geometrically. \marks{2} \\
  \end{enumerate}
\end{enumerate}

\newpage

% ---- Question 7 ----
\questionitem

You are given that $a$ is a parameter which can take only real values. \\
The matrix $A$ is given by
\[
A=\begin{pmatrix}
3&2&a\\
1&3&0\\
2&-1&1
\end{pmatrix}
\] \\
\begin{enumerate}[label=\textbf{(\alph*)}]
  \item Find an expression for the determinant of $A$ in terms of $a$. \marks{2} \\
\end{enumerate}
You are given the following system of equations in $x$, $y$ and $z$.
\[
\begin{aligned}
3x+2y+az&=8\\
x+3y&=5\\
2x-y+z&=3
\end{aligned}
\] \\
The system can be written in the form
\[
A\begin{pmatrix}x\\ y\\ z\end{pmatrix}=\begin{pmatrix}8\\5\\3\end{pmatrix}
\] \\
\begin{enumerate}[label=\textbf{(\alph*)}, resume]
  \item
  \begin{enumerate}[label=(\roman*)]
    \item In the case where $A$ is not singular, solve the given system by using $A^{-1}$. \marks{5}
    \item In the case where $A$ is singular, describe the configuration of the planes whose equations are the three equations of the system. \marks{3} \\
  \end{enumerate}
\end{enumerate}
The transformation represented by $A$ is denoted by $T$. \\
A 3-D solid of volume $a^2+3$ is transformed by $T$ to an image. \\
\begin{enumerate}[label=\textbf{(\alph*)}, resume]
  \item
  \begin{enumerate}[label=(\roman*)]
    \item Determine the range of values of $a$ for which the orientation of the image is the reverse of the orientation of the object. \marks{1}
    \item Determine the range of values of $a$ for which the volume of the image is less than the volume of the object. \marks{2} \\
  \end{enumerate}
\end{enumerate}
\end{enumerate}
\end{document}
